\section{Periodos y torres theta: divisibilidad y normas}
\label{sec:periodos}

\fuente{A, §6, pp.~20--28; B, §4.1, pp.~24--30, y §4.5,
pp.~46--51. Se desarrollan las deducciones lineales, aritméticas y
analíticas que conectan los periodos con la Nota~\ref{sec:marcos}.
La construcción de las torres, las estimaciones locales de jets y las
comparaciones automorfas especializadas se identifican como entradas.}

La divisibilidad de los valores de un emparejamiento permite ampliar
su retículo fuente sin perder la integridad de la aplicación resultante.
Esa ampliación aumenta la pendiente. Para aprovecharla hay que controlar
también la norma compleja del mismo emparejamiento: una ganancia finita
puede quedar compensada por su tamaño arquimediano. Esta nota sigue
ambas contribuciones hasta la cancelación del término principal de la
Nota~\ref{sec:caracteres}.

Los mapas que queremos estudiar tienen la forma
\begin{equation}
 b_\alpha:N_\alpha\longrightarrow C=M_W^\vee,
 \qquad b_\alpha(x)(y)=P_\alpha(x,y),
 \label{eq:periodos-mapa}
\end{equation}
donde $M_W^\vee=\Hom_{\OK}(M_W,\OK)$ lleva la métrica dual.
El cuerpo CM de la construcción se denota por $E$, su subcuerpo
totalmente real por $F$, y el cuerpo de coeficientes por $K$.
No se identifican estos cuerpos. Escribimos $d_K=[K:\Q]$ y usamos
la pendiente duplicada y la altura de
\eqref{eq:marcos-normalizacion}.

\subsection{Los espacios unitarios y las dos presentaciones de las torres}

Un espacio hermitiano sobre $E/F$ es un espacio vectorial sobre $E$
con una forma no degenerada sesquilineal que satisface
$\langle y,x\rangle=\overline{\langle x,y\rangle}$.
Su grupo unitario $U(V)$ conserva esa forma. En cada lugar real de $F$
la forma tiene una signatura $(a,b)$, con $a+b=\dim_E V$.
El grupo real correspondiente es compacto cuando la forma es definida.
El plano hiperbólico $\mathbb H$ tiene una base isotrópica $e,f$ con
$\langle e,f\rangle=1$; su signatura es $(1,1)$.

Sean $k,m\ge2$ pares, y pongamos
\[
 r=k+m-1,\qquad N_j=k+j(t+1),
\]
donde $t\ge1$ es impar. B escribe $h$ en lugar de $t$ y toma
$h\ge r+5$. La siguiente tabla reúne las familias de A, §6, pp.~20--21,
y B, pp.~25--26; los índices recorren las etapas en que se utilizan.

\begin{center}
\small
\begin{tabularx}{\textwidth}{@{}l l Y@{}}
\toprule
Familia & Dimensión sobre $E$ & Datos arquimedianos \\
\midrule
$X_j$ & $N_j$ & Definida en A; positiva definida en la presentación de B. \\
$Y_j$ & $r+jt$ & En B, signatura
$(r+jt-\lfloor(m-j)/2\rfloor,\lfloor(m-j)/2\rfloor)$. \\
$W_j$ & $r$ & En A, signatura $(k+j-1,m-j)$; $W_m$ es definida positiva. \\
$V_j$ & $r+t$ & Espacios auxiliares definidos de la organización de filas de A. \\
\bottomrule
\end{tabularx}
\end{center}

Se tienen inclusiones $X_j\hookrightarrow Y_j$ para $j<m$ y una
inclusión $Y_m\hookrightarrow X_m$ de codimensión uno. En particular,
\[
 \dim_EY_j-\dim_EX_j=m-1-j.
\]
La condición inicial puede escribirse
\[
 Y_0=X_0\mathbin\perp Ez_m\mathbin\perp
                  \mathbb H^{m/2-1},\qquad
 \langle z_m,z_m\rangle=-T_*.
\]
Por ejemplo, $k=m=2$ da $r=3$: las dimensiones de $Y_j$ crecen,
mientras las de $W_j$ permanecen iguales a $3$.

B explica una sucesión de levantamientos desde $X_0$ hasta $X_m$
y otra desde $Y_m$ hacia $Y_0$. A organiza parte del análisis de filas
mediante $W_j$ y $V_j$. Estas son familias distintas: igualar $W_j$
con $Y_j$ sería incompatible con sus dimensiones y signaturas.
La identificación de las realizaciones automorfas mediante las
identidades de periodos es una entrada de la construcción. La tabla
por sí sola no la demuestra.

\subsection{Levantamiento theta y convenciones del emparejamiento}

El anillo de adeles $\mathbb A_F$ reúne las completaciones de $F$:
en los lugares finitos se exige que las coordenadas sean enteras salvo
en un conjunto finito. Escribimos
$[G]=G(F)\backslash G(\mathbb A_F)$ para el cociente automorfo.
Las formas automorfas son funciones sobre ese cociente con las
condiciones de regularidad, crecimiento y finitud de la teoría.
Un levantamiento theta usa un núcleo
\[
 \theta_\varphi(g,h)=\sum_{x\in\mathcal X(F)}
                          (\omega(g,h)\varphi)(x),
\]
donde $\omega$ es la representación de Weil, $\mathcal X$ es un
subespacio lagrangiano de una polarización y $\varphi$ es una función
de Schwartz: en los lugares arquimedianos tiene decaimiento rápido,
y en los finitos se usan funciones localmente constantes de soporte
compacto. La construcción depende de un carácter aditivo, de los
caracteres que fijan el desdoblamiento y de las medidas elegidas.
En la convención lineal en $f$ usada por B,
\begin{equation}
 (\Theta_\varphi f)(h)=c_V\int_{[U(V)]}
                         \theta_\varphi(g,h)f(g)\dd g.
 \label{eq:periodos-theta}
\end{equation}
La convergencia y las propiedades automorfas se reciben aquí de la
teoría theta, con las hipótesis del caso utilizado; véanse B, pp.~24--25,
Gan \cite[§§2--3]{gan-theta} y Gan--Qiu--Takeda
\cite[§11.2]{gqt-rallis}. Escribir una integral formal no justifica
intercambiar su orden ni aplicar una fórmula de norma fuera de su rango.

Para entender la linealidad, el periodo de Bessel de B, p.~27, tiene
el esquema
\begin{equation}
 \mathcal B_\nu(f,g)=c\int_{[U(X_0)]}\int_{[N]}
       \overline{f(h)}\,g(nh)\,\overline{\nu(n)}\dd n\dd h,
 \qquad c>0.
 \label{eq:periodos-bessel-sesqui}
\end{equation}
Aquí $N$ es el grupo unipotente asociado a la bandera elegida y
$\nu$ su carácter. Este periodo es conjugado lineal en $f$ y lineal
en $g$. B realiza el espacio de pruebas mediante funciones conjugadas;
el emparejamiento bilineal correspondiente es
\begin{equation}
 (f,g)\longmapsto\overline{\mathcal B_\nu(f,\overline g)}
   =c\int_{[U(X_0)]}\int_{[N]}f(h)g(nh)\nu(n)\dd n\dd h.
 \label{eq:periodos-bessel-bilineal}
\end{equation}
La segunda igualdad se obtiene conjugando cada factor y usando que
las medidas y $c$ son reales. Las barras forman parte de la convención.
Este prototipo complejo todavía necesita las normalizaciones y las
estructuras sobre $K$ para convertirse en $P_\alpha$.

Las identidades de seesaw comparan los periodos sucesivos con el
emparejamiento final de restricción. La paridad de $m-1-j$ hace alternar
periodos de Bessel y de Fourier--Jacobi; los segundos incorporan otro
factor theta. B, pp.~26--27, atribuye esa comparación a la teoría
clásica y a sus elecciones de datos. La racionalidad del periodo
normalizado y la coincidencia de sus realizaciones son resultados
adicionales, que se conservan como entradas en esta nota.

\subsection{Columnas: sumas finitas y denominadores}

Para un grupo definido a infinito, fijados el peso y el nivel, las
formas algebraicas se describen por vectores de coeficientes sobre un
conjunto finito de clases dobles. B, pp.~27--28, utiliza módulos
integrales de representaciones algebraicas para fijar cuáles de esos
vectores son integrales. En A, §6, p.~21, los operadores de columna adoptan
la forma
\begin{equation}
 (K_j\phi)(L')=\sum_{L\in\mathcal C_j}
       \frac{K_j(L,L')\phi(L)}{|\Gamma_{j,L}|},
 \qquad K_{\rm col}=K_{m-1}\cdots K_0,
 \label{eq:periodos-kernel}
\end{equation}
donde $\Gamma_{j,L}$ es el estabilizador finito de la clase $L$.
Los factores $1/|\Gamma_{j,L}|$ son parte del operador.

Si los coeficientes sin esos pesos preservan los retículos y un entero
$D_\Gamma>0$ es múltiplo de todos los órdenes de los estabilizadores,
cada paso lleva el retículo fuente dentro de $D_\Gamma^{-1}$ veces
el siguiente. Componer $s$ pasos da la inclusión en
$D_\Gamma^{-s}$ veces el retículo final. La prueba consiste en
multiplicar cada suma por $D_\Gamma$ e iterar la inclusión.
Esta observación aísla un denominador; no afirma que sea el único.
Los caracteres de determinante tienen módulo complejo uno, pero esa
propiedad no determina por sí sola sus valoraciones finitas.

Denotemos por $\mathcal U_\alpha$ el espacio de formas de la columna
que B llama $U_\alpha$, con peso y nivel fijados. Se retiene un subespacio
$\mathcal U_\alpha^\circ$ mediante condiciones en el primo auxiliar
$T_\alpha$; aquí se cambia la letra de B para distinguir el espacio
de los unitarios que aparecerán después. Si $\widetilde M_\alpha$
es el retículo inicial, se pone
$M_\alpha=\widetilde M_\alpha\cap\mathcal U_\alpha^\circ$.
El espacio de pruebas común se denota por $\mathcal W$, con norma
de Petersson y retículo $M_W$. Los periodos normalizados son
$K$-bilineales
$P_\alpha:\mathcal U_\alpha^\circ\times\mathcal W\to K$.
El cambio posterior a un subretículo de índice finito
$M'_\alpha\subseteq M_\alpha$ conserva
$\mathcal U_\alpha^\circ$. Esta distinción separa el recorte del
espacio racional del ajuste integral posterior.

\subsection{Filas: el cociente útil y el menor autovalor}

El siguiente lema permite separar el núcleo de una torre de filas sin
perder la información necesaria para acotar la norma del periodo.

\begin{lema}[Descenso por una aplicación sobreyectiva]
Sean $A_0$ y $R_0\ne0$ espacios de Hilbert complejos de dimensión
finita, $T:A_0\to R_0$ lineal y $V\subseteq A_0$ un subespacio.
Supongamos
\[
 T(V)=R_0,\qquad T^*(R_0)\subseteq V.
\]
Definamos $U=V\cap\ker T$ y $Q=V\cap U^\perp$. Entonces
$Q=\operatorname{im}T^*$, y $T|_Q$ es un isomorfismo con inversa
\begin{equation}
 S=T^*(TT^*)^{-1},\qquad
 \|S\|^2=\frac1{\lambda_{\min}(TT^*)}.
 \label{eq:periodos-inversa}
\end{equation}
Si $p:V\to\C$ es lineal y se anula en $U$, existe un único
$\beta:R_0\to\C$ tal que $p=\beta\circ T|_V$, y
\begin{equation}
 \|\beta\|^2\le
       \frac{\|p|_Q\|^2}{\lambda_{\min}(TT^*)}.
 \label{eq:periodos-descenso-norma}
\end{equation}
\label{lem:periodos-cociente}
\end{lema}

\begin{proof}
La sobreyectividad de $T$ implica que $T^*$ es inyectiva y que
$TT^*$ es definida positiva. Además,
$A_0=\ker T\mathbin\oplus^\perp\operatorname{im}T^*$.
La segunda hipótesis permite restringir esta descomposición a $V$:
si $v\in V$, su componente en $\operatorname{im}T^*$ pertenece a
$V$, y la diferencia pertenece a $V\cap\ker T=U$.
Por tanto $V=U\mathbin\oplus^\perp\operatorname{im}T^*$ y
$Q=\operatorname{im}T^*$.

Se tiene $TS=I_{R_0}$ y $S(R_0)\subseteq Q$. Como
$Q\cap\ker T=0$, $S$ es la inversa de $T|_Q$. El cálculo
\[
 S^*S=(TT^*)^{-1}TT^*(TT^*)^{-1}=(TT^*)^{-1}
\]
da \eqref{eq:periodos-inversa} por el teorema espectral.
Si $p|_U=0$, la fórmula $\beta(y)=p(Sy)$ proporciona la
factorización; su unicidad resulta de $T(V)=R_0$.
Finalmente $\|\beta\|\le\|p|_Q\|\,\|S\|$.
\end{proof}

La anulación de $p$ en $U$ es necesaria. La sobreyectividad y una
identidad para $TT^*$ no la implican. Por ejemplo, para
$T:\C^2\to\C$, $T(x,y)=x$, el funcional $p(x,y)=y$ no desciende,
aunque $TT^*=1$. En la construcción formal, esa anulación aparece
como una entrada específica sobre el periodo.

Si una descomposición ortogonal $R_0=\bigoplus_\pi R_0[\pi]$
diagonaliza la torre y
\begin{equation}
 TT^*|_{R_0[\pi]}=
          \kappa_R\prod_{j=1}^m\lambda_{j,\pi}\,I,
 \qquad \kappa_R>0,\quad\lambda_{j,\pi}>0,
 \label{eq:periodos-espectro}
\end{equation}
la norma de la inversa depende del mínimo entre los bloques no nulos.
La igualdad se considera para una tupla de caracteres y una normalización
fijadas: $\kappa_R$ conserva todos los factores positivos comunes a
los bloques, incluidos los que puedan depender de esa tupla.
En B, p.~29, la identidad de la torre $Y_m\to Y_0$ expresa ese
escalar mediante valores $L$ y factores locales. Las fórmulas de
producto interno de Rallis son la entrada especializada correspondiente
\cite[§11]{gqt-rallis}; todavía hay que comprobar sus hipótesis,
los vectores locales y las medidas para los datos de la propuesta.

Aplicando el lema al periodo final $P_\alpha^{\rm fin}$, después
del operador de columnas y con la anulación en $U$ y las
identificaciones de modelos como hipótesis, se obtiene en cada inmersión $\sigma$
la desigualdad de A, (6.1), para el emparejamiento no nulo:
\begin{align}
 \log\|P_\alpha\|_\sigma^2
 &\le\log\|P_\alpha^{\rm fin}\|_\sigma^2
       +\sum_j\log\|K_j\|_\sigma^2-\log\kappa_{R,\sigma}
       \notag\\
 &\hspace{1em}-\min_\pi\sum_j\log\lambda_{j,\pi,\sigma}.
 \label{eq:periodos-cota-arquimediana}
\end{align}
En efecto, se usa $P_\alpha(x,y)=
P_\alpha^{\rm fin}(K_{\rm col}x,Sy)$ y la submultiplicatividad
de la norma. La norma de un emparejamiento bilineal es la de su mapa
al dual. El mínimo espectral se toma antes de promediar sobre
$\sigma$; reemplazarlo por una media de autovalores no controlaría
$\|S\|$.

\subsection{La compatibilidad que necesita el marco}

A, Proposición~6.2, formula la comparación entre el periodo algebraico
y el complejo mediante escalares $a_\alpha\ne0$, unitarios de las
fuentes $U_\alpha$ y un único unitario $V$ del espacio de filas:
\begin{equation}
 P_\alpha(x,y)=a_\alpha\,\operatorname{Bes}_\alpha
                                  (U_\alpha x,Vy).
 \label{eq:periodos-bessel-id}
\end{equation}
El unitario común en el receptor importa cuando se suman operadores.
En términos de mapas hacia un mismo espacio de Hilbert, si
$\widetilde B_\alpha=a_\alpha V B_\alpha U_\alpha$ y
$\widetilde c_\alpha=c_\alpha/|a_\alpha|^2$, entonces
\[
 \widetilde c_\alpha\widetilde B_\alpha^*\widetilde B_\alpha
   =U_\alpha^*(c_\alpha B_\alpha^*B_\alpha)U_\alpha,
 \qquad
 \sum_\alpha\widetilde c_\alpha
       \widetilde B_\alpha\widetilde B_\alpha^*
   =V\left(\sum_\alpha c_\alpha B_\alpha B_\alpha^*\right)V^*.
\]
Así se conservan los defectos del marco de la
Nota~\ref{sec:marcos}. En el paso de un emparejamiento a un mapa
al dual aparece el unitario inducido por $V$, que sigue siendo común.

La comparación no prueba que esos defectos sean pequeños. Tampoco
elimina el escalar de la altura:
\[
 h(a_\alpha B_\alpha)=h(B_\alpha)
       +\frac2{d_K}\sum_\sigma\log|\sigma(a_\alpha)|
\]
cuando $a_\alpha\in K^\times$, $B_\alpha\ne0$ y se mantienen los espacios y las
métricas. Sus valoraciones finitas deben quedar en la misma cuenta.
El cotejo localizado con Lean, descrito más abajo, hace explícita esa
necesidad de transportar a la vez mapa, retículos y métricas.

\subsection{De divisibilidad local a mapa integral}

Sean $M$ y $R$ retículos plenos sobre $\OK$ en espacios
vectoriales sobre $K$, con $\operatorname{rk}M>0$, y sea
$P:M_K\times R_K\to K$ bilineal.
Para cada primo finito $\mathfrak p$ normalizamos
$v_\mathfrak p(\pi_\mathfrak p)=1$. Elegimos enteros
$n_\mathfrak p$, cero salvo en un conjunto finito, y ponemos
$\mathfrak a=\prod_\mathfrak p\mathfrak p^{n_\mathfrak p}$.

\begin{proposicion}[Ampliación por un ideal fraccionario]
Si
\begin{equation}
 P(M_\mathfrak p,R_\mathfrak p)
       \subseteq\pi_\mathfrak p^{n_\mathfrak p}\mathcal O_{K,\mathfrak p}
       \quad\text{para todo }\mathfrak p,
 \label{eq:periodos-divisibilidad}
\end{equation}
entonces $N=\mathfrak a^{-1}M$ satisface $P(N,R)\subseteq\OK$.
Con las mismas métricas arquimedianas,
\begin{equation}
 \mu(N)=\mu(M)+\frac2{d_K}
                 \sum_\mathfrak p n_\mathfrak p\log N\mathfrak p.
 \label{eq:periodos-ganancia-ideal}
\end{equation}
\label{prop:periodos-integridad}
\end{proposicion}

\begin{proof}
Localmente $N_\mathfrak p=\pi_\mathfrak p^{-n_\mathfrak p}
M_\mathfrak p$. La bilinealidad y
\eqref{eq:periodos-divisibilidad} dan valores en
$\mathcal O_{K,\mathfrak p}$. Como
$\OK=K\cap\bigcap_\mathfrak p\mathcal O_{K,\mathfrak p}$,
todos los valores globales son enteros. Si $q=\operatorname{rk}M$,
\[
 \det(\mathfrak a^{-1}M)=\mathfrak a^{-q}\det M.
\]
La fórmula del grado \eqref{eq:marcos-grado} aumenta por
$q\log N\mathfrak a/d_K$. Multiplicar por $2/q$ prueba
\eqref{eq:periodos-ganancia-ideal}.
\end{proof}

La proposición permite exponentes negativos. Si $n_\mathfrak p>0$,
se divide por la potencia correspondiente y aumenta la pendiente;
si es negativo, se contrae el retículo y la pendiente disminuye.
Los ideales fraccionarios no necesitan ser principales.
Esta es la deducción de A, Lema~6.3, y B, p.~49.

Para aplicarla hay que construir antes un único $M_W$ global cuyas
localizaciones sean las utilizadas al medir los periodos. Especificar
retículos locales en todos los primos exige compatibilidad con un
retículo fijo fuera de un conjunto finito. La identificación compleja
\eqref{eq:periodos-bessel-id} no garantiza ese dato integral.

\subsection{El coste del subretículo, el rango y el redondeo}

En la aplicación se parte de $M_\alpha$ y se pasa a un subretículo
$M'_\alpha\subseteq M_\alpha$ de índice finito, del mismo rango
$q_\alpha>0$. Pongamos
\[
 J_{\alpha,\mathfrak p}=\length_{\mathcal O_{K,\mathfrak p}}
       (M_{\alpha,\mathfrak p}/M'_{\alpha,\mathfrak p}),
 \qquad e_\mathfrak p=v_\mathfrak p(p),
\]
donde $p$ es el primo racional bajo $\mathfrak p$. Los exponentes
de la construcción se miden en unidades en que $p$ tiene valoración
uno. Con soporte finito, escribimos
\begin{equation}
 s_{\alpha,\mathfrak p}=\sum_{j=1}^m
     (g_{j,\mathfrak p}+F_{j,\mathfrak p}-E_{\alpha,j,\mathfrak p}),
 \quad n_{\alpha,\mathfrak p}=\lfloor e_\mathfrak p
                                      s_{\alpha,\mathfrak p}\rfloor,
 \quad N_\alpha=\mathfrak a_\alpha^{-1}M'_\alpha.
 \label{eq:periodos-ideal-columna}
\end{equation}
Aquí $\mathfrak a_\alpha=\prod\mathfrak p^{n_{\alpha,\mathfrak p}}$.
La divisibilidad se exige sobre $M'_\alpha\times M_W$.
Combinar \eqref{eq:marcos-coste-indice} y
\eqref{eq:periodos-ganancia-ideal} da la identidad exacta
\begin{equation}
 \mu(N_\alpha)-\mu(M_\alpha)
 =\frac2{d_K}\sum_\mathfrak p
       \left(n_{\alpha,\mathfrak p}
             -\frac{J_{\alpha,\mathfrak p}}{q_\alpha}\right)
                                       \log N\mathfrak p.
 \label{eq:periodos-cuenta-exacta}
\end{equation}
La parte de índice cuesta pendiente, aunque no elimine dimensiones.
El recorte previo del espacio de formas y el paso a un subretículo
de índice finito son operaciones diferentes.

Con $\{x\}=x-\lfloor x\rfloor$, la pérdida exacta por redondeo es
\[
 \rho_\alpha=\frac2{d_K}\sum_\mathfrak p
 \{e_\mathfrak p s_{\alpha,\mathfrak p}\}\log N\mathfrak p\ge0.
\]
Si el soporte no vacío está sobre un conjunto finito $S$ de primos
racionales, entonces
\begin{equation}
 0\le\rho_\alpha<\frac2{d_K}
     \sum_{\mathfrak p\mid S}\log N\mathfrak p
 \le2\sum_{p\in S}\log p.
 \label{eq:periodos-redondeo}
\end{equation}
Si el soporte es vacío, $\rho_\alpha=0$. Para la última desigualdad,
escribimos $N\mathfrak p=p^{f_\mathfrak p}$ y usamos
$\sum_{\mathfrak p\mid p}e_\mathfrak p f_\mathfrak p=d_K$,
de donde $\sum f_\mathfrak p\le d_K$. La cota no crece con el
grado del cuerpo de coeficientes.

Una comparación de rangos sigue siendo necesaria al insertar las cotas
locales del código. Allí una de las longitudes se divide por
$e_\mathfrak p$ y se acota mediante una dimensión $d_\alpha$
anterior al recorte, mientras \eqref{eq:periodos-cuenta-exacta}
divide por el rango retenido $q_\alpha$. Si se dispone de
\[
 \frac{J_{\alpha,\mathfrak p}}{e_\mathfrak p}
       \le d_\alpha(\tau_\mathfrak p+\iota_\mathfrak p),
 \qquad \chi_\alpha=\frac{d_\alpha}{q_\alpha},
\]
la consecuencia correcta es
\begin{equation}
 \frac{J_{\alpha,\mathfrak p}}{e_\mathfrak p q_\alpha}
       -\tau_\mathfrak p
 \le\iota_\mathfrak p+
             (\chi_\alpha-1)(\tau_\mathfrak p+\iota_\mathfrak p).
 \label{eq:periodos-coste-rango}
\end{equation}
La parte positiva del miembro derecho puede incorporarse a una
pérdida no negativa. No hace falta atribuir signos a
$\tau_\mathfrak p$ e $\iota_\mathfrak p$ que la interfaz no declare.
Una cota $q_\alpha\ge(1-\delta_0)d_\alpha>0$ controlaría
$\chi_\alpha\le(1-\delta_0)^{-1}$, pero debe acreditarse.
El primo auxiliar distinguido, excluido de esa cota local, se cuenta
por separado.

\subsection{Por qué aparecen potencias impares de primos}

La parte especializada de la divisibilidad estudia jets de funciones
theta en coordenadas formales y sus módulos integrales. A,
Proposición~6.5, exige hipótesis sobre los logaritmos formales, la
acción de raíces de la unidad, los pesos y la integridad de ciertas
exponenciales. La invariancia por raíces de la unidad, tomada sola,
no da una cota de valoración: multiplicar un monomio invariante por
$p^{-M}$ conserva esa invariancia.

Una identidad elemental explica la forma del término que resulta de
esa estimación. Para $p>1$, $a\ge0$ y números $w_\nu\ge0$,
\begin{equation}
 \frac{p}{p^2-1}a-
       \sum_{\substack{\nu\ge1\\\nu\ {\rm impar}}}
                   \min(a/p^\nu,w_\nu)
 =\sum_{\substack{\nu\ge1\\\nu\ {\rm impar}}}
                   (a/p^\nu-w_\nu)_+.
 \label{eq:periodos-rampa}
\end{equation}
En efecto, $x-\min(x,w)=(x-w)_+$ y
$\sum_{\nu\ge1,\ \nu\ {\rm impar}}p^{-\nu}=p/(p^2-1)$.
Ambas series convergen porque sus términos están entre $0$ y
$a/p^\nu$. Si $w_\nu\ge1$, el miembro derecho solo tiene términos
no nulos cuando $p^\nu<a$. Esta es A, (6.7).

La consecuencia que B, pp.~49--50, utiliza como entrada local es,
con los datos allí construidos,
\begin{equation}
 \sum_{j=1}^m g_{j,\mathfrak p}-\tau_\mathfrak p
 \ge\sum_{j=1}^m\sum_{t\in T_\mathfrak p}
       \sum_{\substack{\nu\ge1\\\nu\ {\rm impar}}}
       \left(\frac{a_t}{p^\nu}
                   -w_{j,c\operatorname{Frob}_p^\nu\tau_t}\right)_+
       -R_\mathfrak p .
 \label{eq:periodos-entrada-local}
\end{equation}
Aquí $T_\mathfrak p$ es el conjunto de lugares seleccionados en los
que se aplica la estimación, $c$ es la conjugación compleja,
$\tau_t$ representa el lugar en el tipo CM y $R_\mathfrak p$ reúne
las pérdidas locales. Los pesos son al menos uno. La fórmula se
aplica en los primos y rangos permitidos por la construcción; no
postula la misma cota en todos los primos.

Para pasar a pendientes, los coeficientes son
\begin{equation}
 c_\mathfrak p=\frac{2e_\mathfrak p\log N\mathfrak p}{d_K},
 \qquad
 \sum_{\mathfrak p\mid p}c_\mathfrak p=2\log p.
 \label{eq:periodos-peso-local}
\end{equation}
La segunda identidad vuelve a usar $\sum e_\mathfrak p
f_\mathfrak p=d_K$. Produce el factor $2\log p$ de la suma
aritmética de la Nota~\ref{sec:caracteres}. La restricción a potencias
impares y el transporte por Frobenius vienen de la estimación local;
la identidad \eqref{eq:periodos-rampa} conserva sus pesos, pero no
prueba esa estimación. Los lugares omitidos de $T_\mathfrak p$ y
los términos $R_\mathfrak p$ necesitan todavía un presupuesto global.

\subsection{La normalización CM del término arquimediano}

Escribamos $f=[F:\Q]$, $\Phi$ para el tipo CM, y
\[
 \operatorname{rd}_E=|\Delta_E|^{1/[E:\Q]},\qquad
 v=\operatorname{rd}_E/2.
\]
En la construcción estudiada, $\kappa$ es el carácter cuadrático
imaginario que intervino en la Nota~\ref{sec:caracteres}, y
$\ell_\kappa=L'(1,\kappa)/L(1,\kappa)$ es real. A, (6.9)--(6.10),
y B, p.~47, fijan
\begin{equation}
 \log\Omega=\gamma+\log(\pi/v)-\ell_\kappa,
 \qquad \Omega>0.
 \label{eq:periodos-cm}
\end{equation}
Aquí $\gamma$ es la constante de Euler. La igualdad procede de
periodos CM y de sus normalizaciones aritméticas; no es una libertad
para ajustar después una constante. A identifica además
$\log\Omega$ con el doble del grado de las líneas CM normalizadas
que usa en los jets. Esa identificación incorpora los datos de
polarización y de comparación de retículos.

Una identidad entre valores absolutos complejos no determina los
denominadores de los factores algebraicos de un periodo. Por eso B,
pp.~46--47, separa la fórmula de periodos, la reciprocidad CM y la
integridad de valores $L$ modificados. Estas comparaciones se reciben
como entradas. Al usarlas habrá que mantener la distinción entre
vectores de Lie y diferenciales duales, y entre norma y norma al
cuadrado: cualquiera de esos cambios altera el signo o el factor
del grado.

\subsection{Cauchy con crecimiento gaussiano}

El origen analítico de las expresiones $a\log a-a$ puede verse en
una estimación completa y elemental. Definimos
\[
 \Phi_0(y)=y\log y-y\quad(y>0),\qquad \Phi_0(0)=0.
\]
Se escribe $\Phi_0$ para distinguir esta función del tipo CM $\Phi$.

\begin{lema}[Derivadas de una función de crecimiento gaussiano]
Sea $G$ entera en $\C^n$ y supongamos
\[
 |G(z)|\le S_0\exp\left(\frac12\sum_{i=1}^n r_i^2|z_i|^2\right),
 \qquad S_0\ge0,\quad r_i>0.
\]
Para un multiíndice $\beta\in\Z_{\ge0}^n$,
\begin{equation}
 |\partial^\beta G(0)|
 \le S_0\prod_{\beta_i>0}
         \beta_i!\left(\frac{e r_i^2}{\beta_i}\right)^{\beta_i/2}.
 \label{eq:periodos-cauchy}
\end{equation}
Si $|c_ir_i|^2=\Omega_i^{-1}$ con $\Omega_i>0$, entonces
\begin{equation}
 |c^\beta\partial^\beta G(0)|^2
 \le S_0^2\prod_i\Omega_i^{-\beta_i}
       \prod_{\beta_i>0} e^2\beta_i\exp(\Phi_0(\beta_i)).
 \label{eq:periodos-cauchy-periodo}
\end{equation}
Los productos vacíos valen uno.
\label{lem:periodos-cauchy}
\end{lema}

\begin{proof}
Podemos fijar en cero las variables para las que $\beta_i=0$.
La fórmula de Cauchy en los círculos $|z_i|=R_i$ da
\[
 |\partial^\beta G(0)|
 \le S_0\prod_{\beta_i>0}
       \beta_i!R_i^{-\beta_i}\exp(r_i^2R_i^2/2).
\]
El logaritmo del factor que depende de $R_i$ tiene derivada
$-\beta_i/R_i+r_i^2R_i$. Su mínimo se alcanza en
$R_i^2=\beta_i/r_i^2$, lo que prueba
\eqref{eq:periodos-cauchy}.

Para un entero $b\ge1$, la concavidad de $\log t$ hace que su
integral en cada intervalo $[j,j+1]$ sea al menos el área del
trapecio de los valores en los extremos. Sumando,
\[
 \log(b!)-\tfrac12\log b
 \le\int_1^b\log t\dd t=\Phi_0(b)+1.
\]
La fórmula también vale para $b=1$. Por tanto
$b!\le e\,b^{b+1/2}e^{-b}$ y
\[
 (b!)^2(e/b)^b\le e^2b\,\exp(\Phi_0(b)).
\]
Elevar \eqref{eq:periodos-cauchy} al cuadrado y multiplicar por
$\prod_i|c_i|^{2\beta_i}$ da
\eqref{eq:periodos-cauchy-periodo}.
\end{proof}

Esta es la estimación de A, Lema~6.6. El factor es
$e^2\beta_i$, con $\beta_i$ multiplicando a $e^2$.
La derivada es ordinaria: sustituirla por el coeficiente de Taylor
$\partial^\beta G(0)/\beta!$ cambiaría la cuenta.

El paso desde derivadas escalares hasta la norma de un bloque de
Schur requiere formas hermitianas y módulos de representaciones
específicos. A, (6.11), obtiene en ese paso factores de la forma
\[
 \Omega_\tau^{-|\nu_\tau|}
 e(2|\nu_\tau|+2)^{2a_0b_0}
 \exp\left(\sum_i\Phi_0(\nu_{\tau,i})\right)
\]
para los bloques de matrices $b_0\times a_0$, bajo las hipótesis de
pesos y relleno por ceros que allí se indican. Se registra como una
entrada especializada. El lema escalar explica el término factorial,
pero no demuestra por sí solo la norma de todo ese bloque.

\subsection{Un cálculo gamma que conserva el término principal}

La fórmula espectral de filas contiene cocientes gamma y factores
de dimensión. El siguiente caso entero muestra cómo extraer su
término principal con una cota uniforme.

\begin{lema}
Sean $a_0\ge\cdots\ge a_{n-1}\ge0$ enteros, $n\ge1$ y $A=a_0$.
Definamos
\[
 H_n(a)=\sum_{i=0}^{n-1}\log\Gamma(a_i+n-i)
       -\sum_{0\le i<l<n}\log(a_i-a_l+l-i).
\]
Entonces
\begin{equation}
 \left|H_n(a)-\sum_{i=0}^{n-1}\Phi_0(a_i)\right|
 \le n+\frac{n(n+1)}2\log(A+n).
 \label{eq:periodos-gamma}
\end{equation}
\label{lem:periodos-gamma}
\end{lema}

\begin{proof}
Los argumentos de los logaritmos son positivos. Como
$\Gamma(a_i+n-i)=a_i!\prod_{h=1}^{n-i-1}(a_i+h)$,
la diferencia entre $H_n(a)$ y $\sum_i\log(a_i!)$ consta de
dos sumas de $n(n-1)/2$ logaritmos, una restada de la otra.
Cada término está entre $0$ y $\log(A+n)$. El valor absoluto
de la diferencia es a lo sumo
$n(n-1)\log(A+n)/2$.

Para $b\ge1$, la monotonía de $\log t$ da
\[
 \Phi_0(b)+1\le\log(b!)
             \le\Phi_0(b)+1+\log b.
\]
Para $b=0$, $\log(0!)=\Phi_0(0)=0$. De aquí
$|\sum_i\log(a_i!)-\sum_i\Phi_0(a_i)|
\le n+n\log(A+n)$. Sumando las dos cotas se obtiene
\eqref{eq:periodos-gamma}.
\end{proof}

El lema es A, Lema~6.7. Las expresiones $\Phi_0$ de las derivadas
y de los factores gamma pueden cancelarse cuando sus pesos y sus
desplazamientos han sido identificados. No se puede concluir esa
cancelación de que ambas estimaciones contengan un término
$a\log a$: los coeficientes, los índices y las constantes lineales
también deben coincidir. A, Proposición~6.8, reúne la comparación
especializada que produce la cota de alturas utilizada a continuación.

\subsection{La cancelación exacta del término principal}

En lo que sigue, con $a_t>0$, $A=\sum_{t\in T}a_t$ y
$w_{j,\rho}>0$, definamos
\begin{align}
 D_{\rm div}
   &=m\sum_{t\in T}\bigl(\Phi_0(a_t)-a_t\log\Omega\bigr),
       \label{eq:periodos-ddiv}\\
 L_w&=\sum_{j=1}^m\frac1f\sum_{\rho\in\Phi}
                           \log(w_{j,\rho}v/\pi),
       \label{eq:periodos-lw}\\
 M_{\rm tot}
   &=\sum_{j=1}^m\sum_{t\in T}a_t
       \left(\log a_t-1-\gamma+\ell_\kappa
                     -\frac1f\sum_{\rho\in\Phi}\log w_{j,\rho}\right).
       \label{eq:periodos-mtot}
\end{align}
El último renglón es la suma de los términos
$\mathcal M(a_t,w_j)$ de \eqref{car-termino-principal}.
Puesto que $|\Phi|=f$, la identidad
\eqref{eq:periodos-cm} implica
\begin{equation}
 \boxed{\quad D_{\rm div}-AL_w=M_{\rm tot}.\quad}
 \label{eq:periodos-cancelacion}
\end{equation}
Para comprobarla, en cada par $(j,t)$ la parte constante es
\[
 -\log\Omega-\frac1f\sum_{\rho\in\Phi}\log(w_{j,\rho}v/\pi)
 =-\gamma+\ell_\kappa-\frac1f\sum_{\rho\in\Phi}\log w_{j,\rho}.
\]
Los términos $\log v-\log\pi$ se cancelan con signos opuestos.
Esta es B, (4.5), p.~48. El cálculo conserva la constante $-1$ de
$\Phi_0(a_t)$ y el signo positivo de $\ell_\kappa$.

\subsection{La desigualdad que enlaza las Notas 4 y 5}

Escribimos $F_{\rm fin}$ y $G_{\rm fin}$ para los números que B
denota por $F$ y $G$ en su p.~49:
\[
 F_{\rm fin}=\sum_\mathfrak p c_\mathfrak p\sum_jF_{j,\mathfrak p},
 \qquad
 G_{\rm fin}=\sum_\mathfrak p c_\mathfrak p
                    \left(\sum_jg_{j,\mathfrak p}-\tau_\mathfrak p\right).
\]
Además, $E_\alpha=\sum_\mathfrak p c_\mathfrak p
\sum_jE_{\alpha,j,\mathfrak p}$.
El resto $\mathcal L_\alpha$ debe cubrir el redondeo, el exceso
del coste de índice sobre la cantidad ya restada mediante
$\tau_\mathfrak p$, y los primos excepcionales. La identidad
\eqref{eq:periodos-cuenta-exacta} indica cómo verificar ese resto;
en particular, no permite perder el cociente de rangos de
\eqref{eq:periodos-coste-rango}.

\begin{proposicion}[Balance condicionado a las estimaciones de periodos]
Supongamos que \eqref{eq:periodos-mapa} es una familia finita no vacía
de mapas integrales, con rangos fuente $q_\alpha>0$ y
$\operatorname{rk}C>0$. Sean $h_\alpha$ números reales con
$h(b_\alpha)\le h_\alpha$ cuando $b_\alpha\ne0$.
Supongamos que se satisfacen las siguientes cotas, con
$E_{\rm div},E_{\rm sl},E_{\rm norm},E_{\rm count}\ge0$
independientes de $\alpha$, y restos $\mathcal L_\alpha\ge0$:
\begin{align}
 \mu(C)+h_\alpha-F_{\rm fin}
   &\le D_{\rm div}-A\log T_\alpha+E_{\rm div},
       \label{eq:periodos-balance-norma}\\
 \mu(N_\alpha)
   &\ge\mu(M_\alpha)+G_{\rm fin}+F_{\rm fin}
                                      -E_\alpha-\mathcal L_\alpha,
       \label{eq:periodos-balance-finito}\\
 \mu(M_\alpha)&\ge-E_{\rm sl},
       \label{eq:periodos-balance-base}\\
 E_\alpha&\le mA\ell_g+E_{\rm norm},
       \label{eq:periodos-balance-normalizacion}\\
 \log T_\alpha&\ge L_w+m\ell_g-E_{\rm count}.
       \label{eq:periodos-balance-conteo}
\end{align}
Aquí $\ell_g=\log(g_0g_1)$ registra dos primos auxiliares y
$T_\alpha$ es el primo de nivel de la columna. Pongamos
\[
 s=\sum_\alpha\frac{q_\alpha}{\sum_\beta q_\beta}
                    \bigl(\mu(N_\alpha)-h_\alpha\bigr),
 \qquad
 \mathcal L=\sum_\alpha\frac{q_\alpha}{\sum_\beta q_\beta}
                    \mathcal L_\alpha.
\]
Entonces
\begin{equation}
 s-\mu(C)\ge G_{\rm fin}-M_{\rm tot}
       -(E_{\rm sl}+E_{\rm norm}+E_{\rm div}
                      +AE_{\rm count})-\mathcal L.
 \label{eq:periodos-balance-media}
\end{equation}
Si se verifican además las hipótesis del
Teorema~\ref{teo:marcos-transferencia}, con esos mapas y esa media,
y si $\mathcal A_{\rm tot}-G_{\rm fin}\le E_{\rm loss}$, se obtiene
\begin{align}
 \mathcal A_{\rm tot}-M_{\rm tot}
 &\le E_{\rm sl}+E_{\rm norm}+E_{\rm div}
                +AE_{\rm count}+\mathcal L+E_{\rm loss}
                \notag\\
 &\quad+\frac{4u\log(1+H_0/u)+(\eta+4\sqrt e)H_0}{1-\delta},
 \label{eq:periodos-exceso}
\end{align}
con la convención del teorema cuando $u=0$.
\label{prop:periodos-balance}
\end{proposicion}

\begin{proof}
Despejar $h_\alpha$ en \eqref{eq:periodos-balance-norma} y
restarlo de \eqref{eq:periodos-balance-finito} cancela
$F_{\rm fin}$ y da
\[
 \mu(N_\alpha)-h_\alpha
 \ge\mu(C)+\mu(M_\alpha)+G_{\rm fin}-D_{\rm div}
       +A\log T_\alpha-E_\alpha
       -E_{\rm div}-\mathcal L_\alpha.
\]
Las tres cotas restantes implican
\[
 A\log T_\alpha-E_\alpha
 \ge AL_w
             -AE_{\rm count}-E_{\rm norm}.
\]
Aquí las dos contribuciones $mA\ell_g$ se
cancelan. Usar \eqref{eq:periodos-cancelacion}, la cota de
$\mu(M_\alpha)$ y promediar con los pesos $q_\alpha/\sum q_\beta$
prueba \eqref{eq:periodos-balance-media}.

El Teorema~\ref{teo:marcos-transferencia} acota $s-\mu(C)$ por
la fracción del segundo renglón de \eqref{eq:periodos-exceso}.
La identidad
\[
 \mathcal A_{\rm tot}-M_{\rm tot}
   =(\mathcal A_{\rm tot}-G_{\rm fin})+(G_{\rm fin}-M_{\rm tot})
\]
y \eqref{eq:periodos-balance-media} prueban la conclusión.
\end{proof}

Las cinco entradas son B, (4.6)--(4.9); la conclusión es (4.10).
$\mathcal A_{\rm tot}$ es la suma de las
$\mathcal A(a_t,t,w_j)$ de \eqref{car-suma-positiva}. Por la
Nota~\ref{sec:caracteres}, su exceso sobre $M_{\rm tot}$ se expresa
en términos de las sumas corregidas $G_\chi$. El denominador
$1-\delta$ afecta a la pérdida del marco; no divide todos los
errores aritméticos.

La proposición demuestra cómo se combinan las estimaciones. La
elección de $T_\alpha$ exige fórmulas de dimensión y masa, junto con
condiciones de descomposición y congruencia; una elección numérica
que solo satisfaga \eqref{eq:periodos-balance-conteo} no realiza
todavía esa construcción. Asimismo, los errores de
\eqref{eq:periodos-exceso} deben ser suficientemente pequeños,
simultáneamente y con sus dependencias uniformes, para producir
las ventanas de la Nota~\ref{sec:ventanas}.

\subsection{Qué queda demostrado y qué debe construirse}

Las deducciones desarrolladas permiten comprobar la cuenta una vez
proporcionados los datos. La aplicación a la propuesta requiere
verificar conjuntamente las entradas siguientes.

\begin{center}
\small
\begin{tabularx}{\textwidth}{@{}>{\raggedright\arraybackslash}p{0.31\textwidth}Y@{}}
\toprule
Entrada & Uso y alcance pendiente \\
\midrule
Realización de las torres y descenso del periodo &
Identificar los modelos sobre $K$, justificar las identidades de
integrales y la anulación en los núcleos. \\
Retículos y divisibilidad local &
Construir $M_W$, $M'_\alpha$ y las cotas de jets, con los índices,
el rango retenido y los primos excepcionales. \\
Normas y periodos CM &
Verificar los escalares de la identidad espectral, las normalizaciones
CM y las estimaciones de bloques que producen
\eqref{eq:periodos-balance-norma}. \\
Geometría, masas y elección del nivel &
Acotar pendientes, dimensiones y $T_\alpha$ para los mismos pesos
y niveles, incluyendo las pérdidas por recortes. \\
Estimación colectiva del marco &
Probar los defectos $e,\eta$ y la proporción retenida del receptor;
las normas individuales no bastan. \\
Presupuesto de errores &
Controlar simultáneamente el miembro derecho de
\eqref{eq:periodos-exceso}, con las dependencias necesarias para
la selección de caracteres y las ventanas. \\
\bottomrule
\end{tabularx}
\end{center}

Las dos primeras operaciones elementales ya tienen pruebas:
descender el periodo por el cociente útil y convertir su divisibilidad
en una aplicación integral. También se ha demostrado el balance
condicional que conecta la identidad de caracteres con la transferencia
de pendientes. El siguiente tramo natural es estudiar las estimaciones
geométricas y de evaluación que suministran las cotas de pendientes y
los defectos colectivos del marco: A, §§7--8, y B, §§4.3--4.4 y~5.

\subsection{Concordancia localizada con Lean}
\label{subsec:periodos-lean}

El cotejo mantiene el espejo de \cite{lean}, revisión
\href{https://github.com/CocoSgt/riemann-hypothesis-lean/tree/7cbf7a0c32fca1badf05931879a37c3f8a2b03b4}{\texttt{7cbf7a0}}.
Se leyeron 19 archivos completos y secciones delimitadas de otros
cuatro. El registro conserva las 23 rutas, URL inmutables, hashes,
declaraciones e hipótesis, y los intervalos de las lecturas selectivas.
Los campos de estructuras y las definiciones de proposiciones se
distinguen de las pruebas que las utilizan.

\begin{enumerate}
\item \textbf{El cociente y la anulación necesaria.}
\href{https://github.com/CocoSgt/riemann-hypothesis-lean/blob/7cbf7a0c32fca1badf05931879a37c3f8a2b03b4/src/Formalism/LocalRingCardinality/DataRows.lean}{\texttt{DataRows}}
define \texttt{V0}, \texttt{Usp} y \texttt{Qsp}, y el adjunto
\texttt{thetaUp} de la composición descendente \texttt{iota}.
En
\href{https://github.com/CocoSgt/riemann-hypothesis-lean/blob/7cbf7a0c32fca1badf05931879a37c3f8a2b03b4/src/Formalism/LocalRingCardinality/DataBu.lean}{\texttt{DataBu}},
\texttt{iotaQ\_inj} prueba la inyectividad en $Q$ y
\texttt{periodIdentity\_iff\_period\_kills}, bajo
\texttt{BuDefQ}, identifica en la tupla de referencia la factorización
del periodo con su anulación en $U$. El enunciado \texttt{FactQZ}
exige la inclusión de la imagen del adjunto en $V$ para concluir
$Q=\operatorname{im}T^*$. Estas declaraciones corresponden
al Lema~\ref{lem:periodos-cociente}; la equivalencia formal no
construye por sí sola la anulación del periodo geométrico.
El espacio superior \texttt{Amb} queda abstracto en
\texttt{DataRows}: no se lo identifica automáticamente con las
formas sobre $Y_m$ de B.

\item \textbf{El mínimo espectral y su normalización.}
La interfaz \texttt{adm1} de
\href{https://github.com/CocoSgt/riemann-hypothesis-lean/blob/7cbf7a0c32fca1badf05931879a37c3f8a2b03b4/src/Formalism/SharpRestriction/Divisibility/LocalQuadraticParts.lean}{\texttt{LocalQuadraticParts}}
emplea $\iota_{\mu,Q}\iota_{\mu,Q}^*$ sobre el espacio inferior
y conserva el factor
\[
 \kappa_R\,\rho(\mu)\prod_j\lambda_{j,\pi}(\mu),
 \qquad \rho(\mu)>0,\quad \rho(\mu_0)=1.
\]
La altura explícita \texttt{hExpl} utiliza la tupla de referencia
$\mu_0$; su ínfimo sobre un conjunto finito no vacío de filas es
el mínimo de \eqref{eq:periodos-cota-arquimediana}, tomado dentro
de cada inmersión. Se leyó \texttt{all\_embeddings\_of\_bounds},
que transporta la cota mediante la métrica de filas, la racionalidad
y la descomposición ortogonal suministradas. No se verificó una
normalización que elimine $\rho(\mu)$ para todas las tuplas.
La escritura de B, p.~29, y esta interfaz deben compararse
conservando ese dato; su diferencia de presentación no se califica
aquí como un error.

\item \textbf{El ideal y el rango de la cota local.}
\href{https://github.com/CocoSgt/riemann-hypothesis-lean/blob/7cbf7a0c32fca1badf05931879a37c3f8a2b03b4/src/Formalism/LocalRingCardinality/DataDIVDefs.lean}{\texttt{DataDIVDefs}}
usa en \texttt{Iu} el exponente
$-\lfloor e_\mathfrak p s_{\alpha,\mathfrak p}\rfloor$:
es el ideal inverso de \eqref{eq:periodos-ideal-columna}.
\texttt{NuDef} identifica, cuando la longitud es finita,
$J_{\alpha,\mathfrak p}=e_\mathfrak p
\operatorname{lenMN}(\alpha,\mathfrak p)$.
Esa igualdad condicional no demuestra por sí sola la finitud.
En los
\href{https://github.com/CocoSgt/riemann-hypothesis-lean/blob/7cbf7a0c32fca1badf05931879a37c3f8a2b03b4/src/Formalism/LocalRingCardinality/RhoDetPrimeEqMuTildeMulChiPow.lean}{datos \texttt{D7}},
\texttt{duS} es la dimensión de \texttt{Mu0}, anterior al corte
\texttt{Mram}. Por eso \texttt{LinkG\_at} lleva al factor
$d_\alpha/q_\alpha$ de \eqref{eq:periodos-coste-rango}.
El primo $T_\alpha$ tiene condiciones separadas.

\item \textbf{El mismo emparejamiento sobre los distintos modelos.}
\href{https://github.com/CocoSgt/riemann-hypothesis-lean/blob/7cbf7a0c32fca1badf05931879a37c3f8a2b03b4/src/Formalism/LocalToGlobal/SiegelLatticeTheory/FlagStabilizerTensors/KerJoinS.lean}{\texttt{KerJoinS.besselId\_of\_joins}}
prueba la comparación unitaria \texttt{adm13} desde las
identificaciones de modelos, el descenso y la sobreyectividad.
El orden de cuantificadores produce un solo $V$ para todas las
columnas. La racionalidad \texttt{BuRational} se añade por separado;
se leyó su prueba de transporte en
\href{https://github.com/CocoSgt/riemann-hypothesis-lean/blob/7cbf7a0c32fca1badf05931879a37c3f8a2b03b4/src/Formalism/LocalToGlobal/Divisibility/DivisibilityRemainder/RationalColumns/BuRationalOf.lean}{\texttt{BuRationalOf}}.
En
\href{https://github.com/CocoSgt/riemann-hypothesis-lean/blob/7cbf7a0c32fca1badf05931879a37c3f8a2b03b4/src/Formalism/LocalToGlobal/Divisibility/RecBridge.lean}{\texttt{RecBridge}},
\texttt{wordFacts\_of\_face} transporta simultáneamente los
retículos y métricas de fuente y filas mediante \texttt{HLatTie},
y el emparejamiento mediante la igualdad \texttt{hBk}.
Invoca \texttt{htMap\_tie} y \texttt{muDual\_tie};
la identificación \texttt{tFo} conserva
$2\,\operatorname{FoOfR}=F_{\rm fin}$.
Estos son puentes demostrados bajo las compatibilidades que reciben,
no una consecuencia de tener dos mapas con nombres parecidos.

\item \textbf{El ensamblaje uniforme y sus entradas.}
Las
\href{https://github.com/CocoSgt/riemann-hypothesis-lean/blob/7cbf7a0c32fca1badf05931879a37c3f8a2b03b4/src/Formalism/SharpRestriction/GlobalTheoremInputs.lean}{interfaces globales}
\texttt{LinkDIV}, \texttt{LinkMu} y \texttt{LinkG} cuantifican
constantes y umbrales comunes. \texttt{LinkDIV\_at} contiene
\[
 s_{\rm frame}+\mu_{\max}(M_W)\le x^{C_H}b
\]
y exige la cota de altura para los mapas no nulos. El factor $b$
de la cota del marco se conserva. La prueba
\texttt{linkDIVt\_of\_rows}, en \texttt{RecBridge}, toma un
máximo de los umbrales proporcionados por \texttt{BridgeRows} y
ensambla las conclusiones. No demuestra desde cero todas las
estimaciones de ese paquete. Del mismo modo,
\texttt{ThetaNormsAndLinks.MetX.I4} ya contiene un criterio de
integridad de jets: las pruebas leídas que lo utilizan no autorizan
a deducir integridad aritmética de una mera norma compleja pequeña.
\end{enumerate}

Los hashes recalculados coinciden con los blobs recuperados en la
revisión fijada. No se compiló ni ejecutó Lean, no se auditó la
clausura de axiomas o dependencias y no se comparó íntegramente el
espejo con la distribución original de Alpöge. La lectura alcanza
las identidades y los puentes especificados; la construcción
automorfa e integral completa, la fórmula de periodos CM y las
cotas uniformes de jets permanecen fuera de ese alcance.

\subsection{Ejercicios con soluciones}

\begin{ejercicio}\label{ej:periodos-minimo}
Sea $T=\operatorname{diag}(\varepsilon,1)$ en $\C^2$, con
$0<\varepsilon<1$. Calcule $\|T^{-1}\|^2$ y la media de los
autovalores de $TT^*$. Explique por qué esa media no puede reemplazar
el mínimo en \eqref{eq:periodos-inversa}.
\end{ejercicio}

\paragraph{Solución del Ejercicio~\ref{ej:periodos-minimo}.}
Los autovalores de $TT^*$ son $\varepsilon^2$ y $1$.
Por tanto $\|T^{-1}\|^2=\varepsilon^{-2}$, mientras su media es
$(1+\varepsilon^2)/2\ge1/2$. El recíproco de la media permanece
menor que $2$, pero $\varepsilon^{-2}$ no está acotado cuando
$\varepsilon$ tiende a cero. Una sola dirección pequeña basta
para aumentar la norma de la inversa.

\begin{ejercicio}\label{ej:periodos-ganancia}
Sobre $\Q$, tome $M=R=\Z^q$, $q\ge1$, con métricas euclídeas y
$P(x,y)=p\sum_i x_iy_i$, donde $p$ es primo.
Construya $N=p^{-1}M$, calcule $\mu(N)$ y la altura de
$b:N\to R^\vee$, y compare $\mu(N)-h(b)$ con $\mu(R^\vee)$.
\end{ejercicio}

\paragraph{Solución del Ejercicio~\ref{ej:periodos-ganancia}.}
La divisibilidad $P(M,R)\subseteq p\Z$ permite dividir el
retículo fuente por $p$. Su covolumen es $p^{-q}$ y
$\mu(N)=2\log p$. El mapa en los espacios euclídeos es
multiplicar por $p$, de modo que su norma es $p$ y
$h(b)=2\log p$. En consecuencia
$\mu(N)-h(b)=0=\mu(R^\vee)$. El ejemplo mantiene las métricas
ambientales y muestra que la ganancia finita puede compensarse
exactamente con la altura.

\begin{ejercicio}\label{ej:periodos-indice}
Sobre $\Q$, sean $M=\Z^2$,
$M'=p\Z\times\Z$ y $N=p^{-1}M'$, con la métrica euclídea.
Compruebe \eqref{eq:periodos-cuenta-exacta} y explique la diferencia
entre índice finito y pérdida de rango.
\end{ejercicio}

\paragraph{Solución del Ejercicio~\ref{ej:periodos-indice}.}
El cociente $M/M'$ tiene longitud local $J_p=1$ y cardinal $p$;
los tres retículos tienen rango $2$. Con $n_p=1$, la fórmula da
$\mu(N)-\mu(M)=2(1-1/2)\log p=\log p$.
Directamente, $N=\Z\times p^{-1}\Z$ tiene covolumen $p^{-1}$,
así que $\mu(N)=\log p$. El índice altera el covolumen sin
disminuir la dimensión del espacio racional.

\begin{ejercicio}\label{ej:periodos-unitario}
Sean $B_1,B_2:\C\to\C^2$ dados por $B_1(z)=ze_1$ y
$B_2(z)=ze_2$, con coeficientes $c_1=c_2=1$. Aplique a la
primera imagen la identidad y a la segunda el unitario que
intercambia $e_1,e_2$. Compare los operadores colectivos antes
y después, aunque las normas individuales permanezcan iguales.
\end{ejercicio}

\paragraph{Solución del Ejercicio~\ref{ej:periodos-unitario}.}
Inicialmente $Q=B_1B_1^*+B_2B_2^*=I_{\C^2}$.
Después ambas imágenes están en $\C e_1$, y
$\widetilde Q=2P_{\C e_1}$. Se tiene
$\operatorname{Tr}Q=\operatorname{Tr}\widetilde Q=2$, pero
$\operatorname{Tr}(Q^2)=2$ y
$\operatorname{Tr}(\widetilde Q^2)=4$.
El defecto $\eta=0$ sirve para la primera familia y la segunda
requiere $\eta\ge1$. Un unitario común habría conservado las
trazas de $Q$ y $Q^2$ mediante conjugación.
